% !TEX program = xelatex
\documentclass[10pt,letterpaper,twocolumn]{article}

\usepackage[top=0.75in, bottom=0.75in, left=0.75in, right=0.75in]{geometry}
\usepackage{fontspec}
\usepackage{unicode-math}
\setmainfont{Latin Modern Roman}
\setsansfont{Latin Modern Sans}
\setmonofont{Latin Modern Mono}[Scale=0.85]
\newfontfamily\acbold{ywft-processing-bold}[Path=../../system/public/type/webfonts/,Extension=.ttf]
\newfontfamily\aclight{ywft-processing-light}[Path=../../system/public/type/webfonts/,Extension=.ttf]

\usepackage{xcolor}
\usepackage{titlesec}
\usepackage{enumitem}
\usepackage{booktabs}
\usepackage{tabularx}
\usepackage{fancyhdr}
\usepackage{hyperref}
\usepackage{graphicx}
\usepackage{ragged2e}
\usepackage{microtype}
\usepackage{natbib}
\usepackage[colorspec=0.92]{draftwatermark}

\definecolor{acpink}{RGB}{180,72,135}
\definecolor{acpurple}{RGB}{120,80,180}
\definecolor{acdark}{RGB}{64,56,74}
\definecolor{acgray}{RGB}{119,119,119}
\definecolor{draftcolor}{RGB}{180,72,135}

\DraftwatermarkOptions{text=WORKING DRAFT,fontsize=3cm,color=draftcolor!18,angle=45,pos={0.5\paperwidth, 0.5\paperheight}}
\hypersetup{colorlinks=true,linkcolor=acpurple,urlcolor=acpurple,citecolor=acpurple,pdfauthor={Jeffrey Alan Scudder},pdftitle={notepat.com: De juguete de teclado a puerta principal del sistema}}

\titleformat{\section}{\normalfont\bfseries\normalsize\uppercase}{\thesection.}{0.5em}{}
\titlespacing{\section}{0pt}{1.2em}{0.3em}
\titleformat{\subsection}{\normalfont\bfseries\small}{\thesubsection}{0.5em}{}
\titlespacing{\subsection}{0pt}{0.8em}{0.2em}

\pagestyle{fancy}\fancyhf{}
\renewcommand{\headrulewidth}{0pt}
\fancyhead[C]{\footnotesize\color{draftcolor}\textit{Borrador de trabajo --- no citar}}
\fancyfoot[C]{\footnotesize\thepage}

\setlist[itemize]{nosep, leftmargin=1.2em, itemsep=0.1em}
\setlist[enumerate]{nosep, leftmargin=1.2em}
\setlength{\columnsep}{1.8em}
\setlength{\parindent}{1em}
\setlength{\parskip}{0.3em}

% Hyphenation for narrow two-column layout
\tolerance=800
\emergencystretch=1em
\hyphenpenalty=50

\newcommand{\acdot}{{\color{acpink}.}}
\newcommand{\ac}{\textsc{Aesthetic Computer}}

\begin{document}

\twocolumn[{%
\begin{center}
{\acbold\fontsize{24pt}{28pt}\selectfont\color{acdark} notepat{\color{acpurple}.}{\color{acpink}com}}\par
\vspace{0.2em}
{\aclight\fontsize{11pt}{13pt}\selectfont\color{acpink} De juguete de teclado a puerta principal del sistema}\par
\vspace{0.6em}
{\normalsize Jeffrey Alan Scudder}\par
{\small\color{acgray} Aesthetic Inc.}\par
{\small\color{acgray} ORCID: \href{https://orcid.org/0009-0007-4460-4913}{0009-0007-4460-4913}}\par
\vspace{0.3em}
{\small\color{acpurple} \url{https://notepat.com} \enspace$\cdot$\enspace \url{https://aesthetic.computer}}\par
\vspace{0.6em}
\rule{\textwidth}{1.5pt}
\vspace{0.5em}
\end{center}

\begin{center}
{\small\color{draftcolor}\textbf{[ borrador de trabajo --- no citar ]}}
\end{center}
\vspace{0.3em}

\begin{quote}
\small\noindent\textbf{Resumen.}
\texttt{notepat} comenzó en junio de 2024 como un sketch de teclado cromático de 200 líneas dentro de \ac{}~\citep{scudder2026ac}. A lo largo de 401 commits y 20 meses creció a 7.903 líneas en el navegador y 987 líneas en metal desnudo, acumulando cinco tipos de forma de onda, reverberación de sala, sensibilidad de presión analógica, un port a Linux sobre metal desnudo que arranca como PID~1, un dominio propio (\texttt{notepat.com}), un arnés de pruebas dedicado, un dispositivo alfa de Ableton Live, un scaffold de plugin VST3, y seis piezas derivadas. Este artículo traza ese crecimiento: en qué empezó \texttt{notepat}, qué funciones ganó, qué infraestructura generó, y cómo una sola pieza llegó a funcionar como la puerta de entrada estratégica de un sistema de computación creativa más grande. Enmarcamos esta trayectoria como una \emph{economía política de una pieza}---los bucles de retroalimentación entre las capacidades crecientes de una pieza, las decisiones de plataforma que motiva y los recursos que atrae a lo largo del tiempo.
\end{quote}
\vspace{0.5em}
}]

\section{Introducción}

Las plataformas de programación creativa se evalúan típicamente a nivel de sistema: diseño de lenguaje, affordances del editor, tamaño de la comunidad~\citep{reas2007processing,mccarthy2015p5js,resnick2009scratch}. Incluso dentro de la investigación de instrumentos musicales digitales, la unidad de análisis suele ser la clase de instrumento o el paradigma de interacción en lugar de la evolución de un solo artefacto a lo largo de años~\citep{magnusson2010designing,marquez2018dmi}. Este artículo toma una lente más estrecha y más larga: una pieza---\texttt{notepat}---dentro de \ac{}, estudiada a lo largo de 20 meses de desarrollo activo.

\texttt{notepat} es un instrumento de teclado melódico. Es directamente accesible como \texttt{aesthetic.computer/notepat} y a través del dominio propio \texttt{notepat.com}. A marzo de 2026, es la pieza más compleja del sistema, la pieza de arranque por defecto en hardware de metal desnudo, y el sujeto principal de la infraestructura de pruebas de rendimiento. Nada de esto fue planeado al inicio. El instrumento creció, y la plataforma creció a su alrededor.

\section{Orígenes (junio--julio 2024)}

El primer commit aparece el 27 de junio de 2024, bajo el nombre \texttt{notepad}---un error tipográfico que fue luego reinterpretado como un portmanteau de ``note'' y ``pad.'' La versión inicial era un muestreador de teclado mínimo: teclas QWERTY mapeadas a notas cromáticas, una sola forma de onda sinusoidal, sin interfaz más allá de rectángulos coloreados para las teclas activas.

En cinco días la pieza había ganado:
\begin{itemize}
  \item Desvanecimiento del sintetizador (decaimiento lineal de 25\,ms para prevenir clics),
  \item Referencias de documentación en el índice del proyecto,
  \item El nombre corregido \texttt{notepat}.
\end{itemize}

En esta etapa la pieza tenía aproximadamente 200 líneas. Usaba el ciclo de vida estándar de pieza de \ac{} (\texttt{boot}, \texttt{paint}, \texttt{act}) y la API integrada \texttt{sound.synth} de la plataforma. El diseño era deliberadamente simple: un sustantivo memorable, un instrumento inmediatamente tocable, sin configuración requerida. Este minimalismo temprano es consistente con lo que \citet{magnusson2010designing} describe como restricción productiva---un espacio de búsqueda acotado que enfoca tanto al intérprete como al desarrollador.

\section{Crecimiento de funciones}

\subsection{Motor de audio (2024--2025)}

La superficie de síntesis se expandió a lo largo de varios ejes impulsados directamente por las necesidades de interpretación. Esta es la dinámica que \citet{mcpherson2020idiomatic} describen en los lenguajes de música por computadora: las herramientas adquieren ``patrones idiomáticos'' que reflejan y refuerzan las preferencias estéticas de sus usuarios principales.

\begin{itemize}
  \item \textbf{Formas de onda}: sinusoidal $\rightarrow$ triángulo, sierra, cuadrada, ruido (seleccionable vía atajo de teclado o parámetro).
  \item \textbf{Reverberación de sala}: línea de retardo de tres taps (intervalos de 120\,ms), coeficiente de retroalimentación 0,55, mezcla húmedo/seco controlable vía deslizador del eje X del trackpad. Suavizado por muestra en metal desnudo ($\alpha \approx 5 \times 10^{-5}$ a 192\,kHz) elimina ruido de cremallera.
  \item \textbf{Cambio de tono}: $\pm$1 octava continua vía eje Y del trackpad, con suavizado exponencial por muestra ($\alpha \approx 3 \times 10^{-4}$ a 192\,kHz).
  \item \textbf{Control de octava}: 9 octavas (1--9), navegables vía teclas de flecha.
  \item \textbf{Modo rápido}: conmutador de ataque/decaimiento rápido para interpretación staccato.
  \item \textbf{Metrónomo}: indicador de pulso sincronizado con tempo (20--300 BPM), destello visual en los tiempos fuertes.
\end{itemize}

\subsection{Entrada y hardware (2024--2026)}

Tres esfuerzos de integración de hardware fueron impulsados por los requisitos de expresividad de \texttt{notepat}. Los instrumentos musicales digitales son descritos por \citet{tahiroglu2020probes} como ``sondas'' que cambian cómo pensamos sobre el sustrato computacional y físico en el que se ejecutan. \texttt{notepat} sondó exactamente de esta manera:

\begin{enumerate}
  \item \textbf{Teclado analógico NuPhy Air60 HE}: Protocolo WebHID con ingeniería inversa de capturas USB. Comandos de activación de sesenta y cuatro bytes; reportes de presión de 16 bits por tecla; síntesis ponderada por velocidad. Requirió \texttt{CONFIG\_HIDRAW=y} en el kernel de metal desnudo.
  \item \textbf{Trackpad como controlador de expresión}: El eje X controla la mezcla húmeda de reverberación, el eje Y controla el cambio de tono. Estos mapeos fueron diseñados para el estilo de interpretación de \texttt{notepat} y no se han generalizado.
  \item \textbf{Pantalla HDMI secundaria}: Salida de color sólido impulsada por notas mezcladas durante la interpretación. Añadida a \texttt{ac-native} en marzo de 2026, motivada por uso en vivo.
\end{enumerate}

\subsection{Diseño visual}

La interfaz evolucionó de rectángulos coloreados simples a un sistema en capas:

\begin{itemize}
  \item Cuadrícula dividida de 24 pads (4$\times$3 octava izquierda + 4$\times$3 octava derecha),
  \item Barra de estado: reloj (hora LA con horario de verano), forma de onda, octava, tasa de muestreo, volumen, batería, icono de antena WiFi,
  \item Visualización de mini teclado de piano con sombreado de octava y notación de puntos para teclas negras,
  \item Etiqueta HUD en fuente MatrixChunky8 con marca de \texttt{.com} en superíndice,
  \item Tinte de fondo reactivo a la hora del día (amanecer/mediodía/tarde/atardecer/noche),
  \item Sistema de color de notas: cada semitono mapeado a un matiz (Do = rojo hasta Si = violeta), las notas activas se mezclan en el fondo.
\end{itemize}

\subsection{Integración con DAW: Ableton Live}

En enero de 2025, se creó un dispositivo alfa de Max for Live (\texttt{AC Notepat.amxd}) para incrustar \texttt{notepat} dentro de Ableton Live. El dispositivo usa el objeto \texttt{jweb\textasciitilde} de Max~9---un navegador web incrustado con audio habilitado---para cargar la pieza \texttt{notepat} y enrutar su salida Web Audio directamente al mezclador de Ableton vía \texttt{plugout\textasciitilde}.

Esta integración requirió resolver varios problemas no presentes en los contextos de navegador o metal desnudo:

\begin{itemize}
  \item \textbf{Empaquetado offline}: Un generador de bundles HTML autocontenidos (\texttt{bundle-html.js}) descubre recursivamente todas las dependencias de módulos ES, inserta glifos de fuentes en línea, y produce un archivo HTML auto-extraíble comprimido con gzip con un sistema de archivos virtual (VFS) e interceptor de fetch.
  \item \textbf{AudioWorklet en modo PACK}: La ruta estándar de carga de AudioWorklet requiere un servidor; la versión empaquetada crea URLs blob desde el VFS para cargar un worklet de altavoz pre-empaquetado (\texttt{speaker-bundled.mjs}, $\sim$50\,KB).
  \item \textbf{Traducción MIDI-a-teclado}: Una capa de mapeo planeada convierte eventos MIDI note-on/off al modelo de eventos de teclado de \texttt{notepat}, con CC1 (rueda de modulación) controlando la mezcla de reverberación y pitch bend controlando el modo de deslizamiento.
\end{itemize}

Un esfuerzo paralelo de plugin VST3 (\texttt{ac-vst/}) armó un scaffold de plugin C++ con SDK de Steinberg con una UI basada en WebView, apuntando a una ruta de distribución más portable que Max for Live. Ambos esfuerzos permanecen en alfa; el motor de audio aún no es funcional en el contexto empaquetado offline debido a restricciones de carga de worklets.

La integración con DAW representa un quinto bucle de retroalimentación en la economía política de \texttt{notepat}: \textbf{presión de distribución $\rightarrow$ infraestructura de empaquetado}. El deseo de hacer \texttt{notepat} disponible dentro de entornos de producción musical profesional forzó a la plataforma a desarrollar capacidades de empaquetado HTML offline, emulación de VFS y pre-empaquetado de worklets que ninguna otra pieza había requerido.

\subsection{Integración de red y sistema}

\begin{itemize}
  \item Interfaz WiFi: escáner de red a pantalla completa con barras de intensidad de señal, selector de SSID, entrada de contraseña---construido en la propia pieza para contextos de metal desnudo donde no existe interfaz de SO.
  \item Modo oscuro: cambio automático a las 7\,pm/7\,am hora de LA, con cálculo de horario de verano correspondiente en tanto C como JavaScript.
  \item Modos de proyector y recital para contextos de actuación.
\end{itemize}

\section{El port a metal desnudo}

A finales de 2024, \texttt{notepat} fue portado a \texttt{ac-native}: un entorno de ejecución Linux personalizado que arranca desde USB como PID~1 sin sistema operativo. La versión de metal desnudo (987 líneas) se ejecuta en QuickJS, renderiza vía búferes tontos DRM/KMS con escalado CPU 3$\times$ por vecino más cercano, y sintetiza audio a través de acceso directo ALSA \texttt{hw:} a 192\,kHz.

Incrustar un sintetizador de audio en tiempo real en un entorno Linux de metal desnudo es un área activa de investigación. \citet{turchet2021elk} describen Elk Audio OS---un SO basado en Linux que usa extensiones de kernel en tiempo real Xenomai para latencia de audio sub-milisegundo en hardware embebido---como el estado del arte en este espacio. El enfoque de \texttt{ac-native} difiere: en lugar de un SO de audio embebido de propósito general, es un entorno de ejecución de pieza única sin capa RTOS, confiando en lugar de ello en búferes de periodo ALSA grandes (128 muestras a 192\,kHz) y un motor de síntesis solo-software con polifonía de 32 voces.

Este port no fue incidental. Todo el proyecto \texttt{ac-native}---configuración de kernel personalizada, pipeline de compilación de initramfs, motor de audio ALSA, manejo de entrada evdev, gestión de WiFi vía \texttt{wpa\_supplicant}, pantalla DRM/KMS---fue construido principalmente para ejecutar \texttt{notepat} en hardware físico como un instrumento dedicado. La pieza justificó la plataforma.

La secuencia de arranque refleja este compromiso: después del init del kernel, \texttt{ac-native} muestra un título animado en arcoíris, luego carga \texttt{notepat.mjs} como la pieza por defecto. El instrumento es tocable aproximadamente 2 segundos después del encendido.

\section{Derivados y reutilización}

Seis artefactos derivan directamente de \texttt{notepat}:

\begin{enumerate}
  \item \textbf{\texttt{autopat.mjs}} (214 líneas): Wrapper de reproducción automática con un jukebox integrado (Twinkle, extracto de Bach, escalas). Importa y delega a las funciones centrales del ciclo de vida de \texttt{notepat}.
  \item \textbf{\texttt{notepat-tv.mjs}} (55 líneas): Visualizador remoto dirigido por UDP para instalaciones artísticas. Recibe datos de notas y dirige gráficos de tortuga en una superficie proyectada.
  \item \textbf{\texttt{notepat-convert.mjs}} (103 líneas): Biblioteca de traducción de notación; convierte atajos de teclado de una sola letra a sintaxis de melodía extendida. Reutilizada por \texttt{clock}, \texttt{stample} y otras piezas musicales.
  \item \textbf{\texttt{notepat.mjs} de metal desnudo} (987 líneas): Reimplementación independiente para QuickJS, compartiendo mapeos de teclado y convenciones de interfaz pero adaptada para acceso directo al hardware.
  \item \textbf{\texttt{AC Notepat.amxd}} (dispositivo Max for Live): Incrusta \texttt{notepat} dentro de Ableton Live vía \texttt{jweb\textasciitilde}, enrutando la síntesis Web Audio directamente al mezclador del DAW. Alfa, enero 2025.
  \item \textbf{Plugin \texttt{ac-vst}} (C++/SDK Steinberg): Scaffold VST3 con UI WebView para distribución agnóstica de DAW. Incluye manejo de entrada MIDI, enrutamiento de buses de audio y stubs de automatización de parámetros.
\end{enumerate}

\section{Infraestructura generada}

\subsection{Enrutamiento y marca}

El dominio \texttt{notepat.com} se reescribe en el edge hacia el slug \texttt{notepat} en la función de índice principal. Esto requirió lógica dedicada de coincidencia de host, metadatos Open Graph personalizados, y un elemento HUD de marca visible durante la reproducción. \texttt{notepat} es la única pieza en el sistema con su propio dominio. El superíndice \texttt{.com} visible en el HUD del instrumento codifica esta doble identidad en la superficie de actuación misma.

\subsection{Infraestructura de pruebas}

\texttt{notepat} tiene una inversión de pruebas desproporcionada en relación con otras piezas:

\begin{itemize}
  \item \textbf{Arnés de latencia} (\texttt{artery/test-notepat-latency.mjs}): Mide la latencia de teclado-a-sonido vía Chrome DevTools Protocol; reporta media, mediana, p95, p99. Medición actual: 417\,ms (objetivo: $<$30\,ms).
  \item \textbf{Arnés de estabilidad} (\texttt{artery/test-notepat-stability.mjs}): Prueba de fuzz de 5 minutos a intensidad configurable (2--20 notas/seg); monitorea crecimiento del heap, tamaño del diccionario de sonido, degradación de latencia. Último resultado: 3.229 notas, $-$14\% heap, 0 advertencias.
  \item 14 scripts de prueba en total a través de \texttt{artery/}, \texttt{tests/} y \texttt{.vscode/tests/}.
\end{itemize}

La brecha de latencia (417\,ms vs.\ 30\,ms objetivo) representa un problema de investigación abierto. Es consistente con las restricciones conocidas de programación de audio del navegador y el costo del despacho de eventos de la API Web Audio desde JavaScript. Esta brecha no surge en la versión de metal desnudo, donde el acceso ALSA \texttt{hw:} a 192\,kHz y el sondeo evdev directo producen latencia de onset sub-10\,ms.

\subsection{Analíticas de hits de piezas}

El sistema de seguimiento de hits de piezas (\texttt{/api/piece-hit}) fue construido en diciembre de 2025. Registra cargas de página del lado del servidor en MongoDB, distinguiendo piezas del sistema y de usuario. En la ventana de telemetría inicial (2025-12-31 a 2026-01-02), \texttt{notepat} registró 27 hits y se clasificó 4to entre 44 piezas integradas rastreadas. Las limitaciones de instrumentación actuales---sin espacios de nombres tipados, sin atribución autenticada, llamadas POST de disparar-y-olvidar que estaban siendo terminadas por el runtime de Netlify Functions antes de completarse---significan que la imagen completa del uso de \texttt{notepat} desde 2024 no está capturada. El error de disparar-y-olvidar fue identificado y corregido en marzo de 2026.

\section{Economía política de una pieza}

La ``ley de rendimientos acelerados'' de Kurzweil describe cómo los sistemas tecnológicos acumulan capacidad a una tasa acelerada a medida que cada avance habilita el siguiente~\citep{kurzweil2005singularity}. A la escala de una sola pieza creativa, una dinámica similar es visible. La trayectoria de \texttt{notepat} ilustra bucles de retroalimentación entre una pieza y su plataforma anfitriona que \citet{nieborg2018platformization} describen a nivel de industrias culturales: las affordances de la plataforma moldean la producción, y las demandas de producción remoldean las plataformas.

\begin{enumerate}
  \item \textbf{Presión de funciones $\rightarrow$ capacidad de plataforma}: \texttt{notepat} necesitaba reverberación de sala, así que el motor de audio ganó líneas de retardo. Necesitaba presión analógica, así que \texttt{ac-native} ganó soporte HID raw. Necesitaba una pantalla secundaria, así que se añadió salida HDMI al backend DRM.
  \item \textbf{Presión de marca $\rightarrow$ decisiones de enrutamiento}: El deseo de una URL memorable llevó a \texttt{notepat.com}, que requirió reescritura de funciones edge de host y vistas previas sociales personalizadas.
  \item \textbf{Presión de calidad $\rightarrow$ inversión en pruebas}: Una regresión de latencia en \texttt{notepat} motivó la construcción de un arnés de pruebas de piezas de propósito general (artery), que ahora beneficia a todas las piezas.
  \item \textbf{Presión de hardware $\rightarrow$ entorno de ejecución completo}: La ambición de ejecutar \texttt{notepat} en hardware dedicado justificó construir \texttt{ac-native}: kernel personalizado, initramfs, motor ALSA, pantalla DRM, pila WiFi---miles de líneas de C.
  \item \textbf{Presión de distribución $\rightarrow$ infraestructura de empaquetado}: El deseo de incrustar \texttt{notepat} dentro de Ableton Live forzó a la plataforma a desarrollar empaquetado HTML offline con un sistema de archivos virtual, interceptor de fetch y AudioWorklets pre-empaquetados---infraestructura que ninguna otra pieza había requerido.
\end{enumerate}

\citet{born2015music} observan que la inversión en tecnología musical nunca es neutral---refleja y refuerza las prioridades estéticas, posiciones sociales y recursos de quienes la construyen y financian. La trayectoria de \texttt{notepat} hace esto visible a la escala de un desarrollador individual y un único instrumento de código abierto: la inversión personal sostenida en una pieza moldea materialmente lo que es técnicamente posible para cada otra pieza del sistema.

Las comunidades de programación en vivo han desarrollado un vocabulario explícito para este tipo de relación entre práctica y herramienta~\citep{blackwell2022livecoding}: el instrumento y el instrumentista co-evolucionan. Los 401 commits de \texttt{notepat} son un registro de esa co-evolución en un contexto no de programación en vivo---un instrumento persistente que acumuló capacidad a través del uso en lugar del diseño previo.

\section{Escala actual}

La Tabla~\ref{tab:scale} resume la huella de \texttt{notepat} al 8 de marzo de 2026.

\begin{table}[h]
\small
\centering
\begin{tabularx}{\columnwidth}{l r}
\toprule
\textbf{Métrica} & \textbf{Valor} \\
\midrule
Commits (versión navegador) & 401 \\
LOC navegador & 7.903 \\
LOC metal desnudo & 987 \\
Piezas derivadas & 6 \\
LOC total del ecosistema & 9.262 \\
Tipos de forma de onda & 5 \\
Scripts de prueba dedicados & 14 \\
Archivos que mencionan ``notepat'' en el repo & 159 \\
Edad (primer commit) & 27 de junio de 2024 \\
\bottomrule
\end{tabularx}
\caption{Indicadores de escala de notepat (8 de marzo de 2026).}
\label{tab:scale}
\end{table}

\section{Conclusión}

\texttt{notepat} demuestra que una sola pieza puede funcionar como el centro de gravedad estratégico de una plataforma de computación creativa. Su crecimiento de un sketch de teclado de 200 líneas a un ecosistema de 9.000 líneas con su propio dominio, entorno de ejecución de hardware, integración con DAW e infraestructura de pruebas no fue planeado---emergió del uso sostenido y los bucles de retroalimentación descritos arriba.

La contribución metodológica es una unidad de análisis: la \emph{trayectoria de pieza}. Estudiar una pieza individual a lo largo de su historial completo de commits---qué funciones ganó, qué infraestructura generó, qué derivados engendró---revela dinámicas de plataforma invisibles a nivel de sistema. La investigación de instrumentos musicales digitales ha examinado la longevidad de los DMI como un problema de adopción~\citep{marquez2018dmi} y los DMIs como sondas del sustrato computacional~\citep{tahiroglu2020probes}. La trayectoria de pieza combina ambas perspectivas con una lente de economía política: los instrumentos acumulan capacidad, y esa acumulación tiene consecuencias materiales para los sistemas que los albergan.

\vspace{0.5em}
\noindent\textit{Traducido del inglés. Versión original disponible en \url{https://papers.aesthetic.computer}}

\bibliographystyle{plainnat}
\bibliography{references}

\end{document}
